\item
@@PAGE_MARK@@For every extension @@M0@@ of @@M1@@, the kernel of @@M2@@ is a torsion group.
\item For every pure extension @@M3@@ of @@M4@@, the homomorphism @@M5@@ is injective. (NB. Note that @@M6@@ is not necessarily assumed infinite).
\end{enumerate}

In case (i), one is reduced to proving that every element @@M7@@ in the kernel of @@M8@@, where @@M9@@ is an integral algebra of finite type over @@M10@@, is a torsion element. Using a @@M11@@-homomorphism @@M12@@, where @@M13@@ is a finite @@M14@@-algebra, one is reduced to the case where @@M15@@ is a finite extension of @@M16@@. One then has @@M17@@ by the property postulated above for the norm homomorphism. This proves that @@M18@@ is a torsion element. One even has @@M19@@ when the greatest common divisor of the degrees of the finite extensions @@M20@@ of @@M21@@ that are quotients of @@M22@@ equals 1. This is the case when the fraction field @@M23@@ of @@M24@@ is a pure extension of @@M25@@, as is readily checked (only the case where @@M26@@ is finite requires a new verification). This observation then gives the conclusion of (ii).

\sgaStableTarget{sga6.e10.para.7-18-4}\textbf{7.18.4.} Suppose that, for every nonempty affine open subset @@M27@@ of the affine space @@M28@@, with @@M29@@ the ring of @@M30@@, the homomorphism @@M31@@ is surjective (compare IX 1.6). Then, using the commutation of @@M32@@ with filtered inductive limits, one finds that for every pure extension @@M33@@ of @@M34@@, @@M35@@ is surjective. If one also assumes that @@M36@@ is infinite, or that norm homomorphisms can be defined as in 7.18.3, one concludes that this homomorphism is even bijective.

\sgaStableTarget{sga6.e10.para.7-18-5}\textbf{7.18.5.} As an example, consider a scheme @@M37@@ of finite type over @@M38@@, and set
@@M39@@
